\chapter{Analisi dei risultati}
\label{ch:analisi}

In questa sezione viene fornita un'analisi dei dati derivati dagli esperimenti condotti con gli utenti. Tale analisi si articola in tre macro-aree: i profili demografici estratti dal questionario iniziale, l'analisi puntuale compito per compito, e i risultati dei questionari finali (SUS e NPS) corredati dai commenti qualitativi.

\section{Profili demografici sperimentatori}

Un totale di 12 utenti è stato selezionato e diviso in due gruppi (A e B), ciascuno composto da 6 partecipanti. Di seguito vengono riportati i risultati ottenuti dall'analisi anagrafica, necessaria per garantire che i campioni fossero bilanciati e rappresentativi. Non essendovi sbilanciamenti severi si è assunto che i gruppi fossero idonei all'esperimento \textit{within-subject}.

\subsection{Genere}
Il campione è composto da 6 uomini e 6 donne, perfettamente bilanciato. Nel Gruppo A la distribuzione è di 3 uomini (U01, U05, U06) e 3 donne (U02, U03, U04); nel Gruppo B si hanno parimenti 3 uomini (U07, U09, U12) e 3 donne (U08, U10, U11). La distribuzione è del tutto equa tra i gruppi.

\subsection{Età}
L'età dei partecipanti spazia dai 19 ai 57 anni, con due cluster demografici netti e intenzionali:
\begin{itemize}
    \item \textbf{Cluster giovani} (19--22 anni): U04, U05, U06, U10, U11, U12 — alta familiarità tecnologica, esperienza con prenotazioni online.
    \item \textbf{Cluster maturi} (51--57 anni): U01, U02, U03, U07, U08, U09 — competenze informatiche più limitate, minore abitudine alle prenotazioni digitali.
\end{itemize}
Questa bipartizione è stata deliberata per testare l'accessibilità trasversale della piattaforma, verificando se V2 risulta migliorata sia per utenti esperti che per utenti meno avvezzi alla tecnologia.

\subsection{Titolo di studio e Competenze Informatiche}
Il campione presenta una distribuzione variegata: 3 partecipanti con Laurea Triennale (U04, U10, U12), 5 con Diploma (U03, U05, U06, U08, U11) e 4 con Licenza Media (U01, U02, U07, U09). Le competenze informatiche autovalutate su scala 1--5 riflettono fedelmente la bipartizione anagrafica: il cluster giovane si attesta su valori 3--5 (Media/Alta), mentre il cluster maturo registra valori 1--2 (Bassa). Entrambi i gruppi A e B includono una rappresentanza bilanciata di utenti ad alta e bassa competenza.

\subsection{Autonomia nelle prenotazioni online}
Il questionario anagrafico includeva una domanda sull'autonomia nell'effettuare prenotazioni online (scala 0--2: mai / a volte / sempre). Il dato è particolarmente rilevante per una piattaforma di prenotazione visite: 7 partecipanti su 12 dichiarano piena autonomia (valore 2), 4 autonomia parziale (valore 1) e 1 partecipante (U01) non ha mai effettuato prenotazioni online in autonomia. Questo garantisce la presenza di utenti "nuovi" al paradigma di prenotazione digitale, rendendo il test più significativo.

\subsection{Daltonismo}
Nessun partecipante ha dichiarato di soffrire di daltonismo, tuttavia, a livello precauzionale e di standard di accessibilità, l'interfaccia V2 mantiene un rapporto di contrasto a norma WCAG e affida i feedback visivi non solo al colore ma anche a etichette esplicite (es. Status Badge azzurro per "Upcoming", grigio per "Complete").

\section{Analisi dei compiti}

In questa sezione viene esposta l'analisi dei dati raccolti. Per ciascun compito (da 1 a 11) i dati quantitativi sono stati elaborati confrontando le metriche ottenute nella Versione 1 (originale) e nella Versione 2 (riprogettata). Indichiamo con "C" i compiti completati in autonomia, "CA" per completati con aiuto e "NC" per non completati.

\subsection{Compito 1: Sede organizzazione}
\begin{table}[H]
\centering
\renewcommand{\arraystretch}{1.2}
\small
\begin{tabular}{|l|c|c|}
\hline
\textbf{Metrica} & \textbf{Versione 1} & \textbf{Versione 2} \\
\hline
Completato & 8 (C), 1 (CA), 1 (NC) & 10 (C), 0 (CA), 0 (NC) \\ \hline
Tempo Medio & 75 s & 22 s \\ \hline
Click Medi & 6 & 2 \\ \hline
Errori Medi & 1.2 & 0.1 \\ \hline
Difficoltà Percepita & Facile & Molto facile \\ \hline
\end{tabular}
\end{table}
\textbf{Commenti:} Nella V1, l'indirizzo era nascosto all'interno di sottomenu testuali; nella V2 è immediatamente visibile nel footer universale. Tempi di completamento abbattuti di oltre il 60\%.

\subsection{Compito 2: Registrazione e accesso}
\begin{table}[H]
\centering
\renewcommand{\arraystretch}{1.2}
\small
\begin{tabular}{|l|c|c|}
\hline
\textbf{Metrica} & \textbf{Versione 1} & \textbf{Versione 2} \\
\hline
Completato & 7 (C), 2 (CA), 1 (NC) & 10 (C), 0 (CA), 0 (NC) \\ \hline
Tempo Medio & 148 s & 92 s \\ \hline
Click Medi & 13 & 8 \\ \hline
Errori Medi & 1.9 & 0.3 \\ \hline
Difficoltà Percepita & Media & Facile \\ \hline
\end{tabular}
\end{table}
\textbf{Commenti:} L'aggiunta del toggle visivo per la password e la corretta gestione degli input di form in V2 hanno azzerato gli errori di battitura fatali durante la creazione dell'account.

\subsection{Compito 3: Prenotazione a "Castello di Sera"}
\begin{table}[H]
\centering
\renewcommand{\arraystretch}{1.2}
\small
\begin{tabular}{|l|c|c|}
\hline
\textbf{Metrica} & \textbf{Versione 1} & \textbf{Versione 2} \\
\hline
Completato & 6 (C), 3 (CA), 1 (NC) & 10 (C), 0 (CA), 0 (NC) \\ \hline
Tempo Medio & 185 s & 110 s \\ \hline
Click Medi & 15 & 7 \\ \hline
Errori Medi & 2.5 & 0.4 \\ \hline
Difficoltà Percepita & Media/Difficile & Facile \\ \hline
\end{tabular}
\end{table}
\textbf{Commenti:} In V1 l'assenza di filtro (\textit{Faceted Search}) ha disorientato gli utenti. La nuova griglia a card con status espliciti della V2 ha facilitato incredibilmente l'individuazione di questo evento specifico.

\subsection{Compito 4: Codice prenotazione}
\begin{table}[H]
\centering
\renewcommand{\arraystretch}{1.2}
\small
\begin{tabular}{|l|c|c|}
\hline
\textbf{Metrica} & \textbf{Versione 1} & \textbf{Versione 2} \\
\hline
Completato & 8 (C), 1 (CA), 1 (NC) & 10 (C), 0 (CA), 0 (NC) \\ \hline
Tempo Medio & 68 s & 31 s \\ \hline
Click Medi & 6 & 2 \\ \hline
Errori Medi & 1.4 & 0.1 \\ \hline
Difficoltà Percepita & Media & Molto facile \\ \hline
\end{tabular}
\end{table}
\textbf{Commenti:} La V2 vanta un layout ``stile biglietto'' per le prenotazioni. L'\textit{affordance} grafica garantita dai font monospazio e dal grasettaggio del token ha reso l'informazione saliente al primo sguardo.

\subsection{Compito 5: Cancella prenotazione}
\begin{table}[H]
\centering
\renewcommand{\arraystretch}{1.2}
\small
\begin{tabular}{|l|c|c|}
\hline
\textbf{Metrica} & \textbf{Versione 1} & \textbf{Versione 2} \\
\hline
Completato & 5 (C), 3 (CA), 2 (NC) & 10 (C), 0 (CA), 0 (NC) \\ \hline
Tempo Medio & 175 s & 48 s \\ \hline
Click Medi & 12 & 3 \\ \hline
Errori Medi & 2.8 & 0.2 \\ \hline
Difficoltà Percepita & Difficile & Facile \\ \hline
\end{tabular}
\end{table}
\textbf{Commenti:} Uno dei task peggiori per V1. Il pulsante fuorviante confondeva gli utenti. In V2, l'introduzione della modale rossa e dell'invito standardizzato non hanno richiesto alcun intervento da parte del facilitatore.

\subsection{Compito 6: Switch account volontario}
\begin{table}[H]
\centering
\renewcommand{\arraystretch}{1.2}
\small
\begin{tabular}{|l|c|c|}
\hline
\textbf{Metrica} & \textbf{Versione 1} & \textbf{Versione 2} \\
\hline
Completato & 9 (C), 1 (CA), 0 (NC) & 10 (C), 0 (CA), 0 (NC) \\ \hline
Tempo Medio & 45 s & 35 s \\ \hline
Click Medi & 5 & 4 \\ \hline
Errori Medi & 0.4 & 0.1 \\ \hline
Difficoltà Percepita & Facile & Molto facile \\ \hline
\end{tabular}
\end{table}
\textbf{Commenti:} Entrambe le versioni hanno performato bene in quanto l'azione di logout e re-login segue \textit{pattern de facto} standard (corner in alto a destra). Leggero vantaggio di V2 dovuto ai campi pre-formattati per l'email.

\subsection{Compito 7: Imposta disponibilità}
\begin{table}[H]
\centering
\renewcommand{\arraystretch}{1.2}
\small
\begin{tabular}{|l|c|c|}
\hline
\textbf{Metrica} & \textbf{Versione 1} & \textbf{Versione 2} \\
\hline
Completato & 6 (C), 3 (CA), 1 (NC) & 10 (C), 0 (CA), 0 (NC) \\ \hline
Tempo Medio & 155 s & 52 s \\ \hline
Click Medi & 10 & 4 \\ \hline
Errori Medi & 1.8 & 0.2 \\ \hline
Difficoltà Percepita & Medio/Difficile & Facile \\ \hline
\end{tabular}
\end{table}
\textbf{Commenti:} L'immissione della data manuale su input testuale in V1 (suscettibile a errori per \textit{date format}) è stata soppiantata da un widget calendario (Datepicker) nativo in V2. Gli errori sono magicamente collassati.

\subsection{Compito 8: Switch account admin}
\begin{table}[H]
\centering
\renewcommand{\arraystretch}{1.2}
\small
\begin{tabular}{|l|c|c|}
\hline
\textbf{Metrica} & \textbf{Versione 1} & \textbf{Versione 2} \\
\hline
Completato & 10 (C), 0 (CA), 0 (NC) & 10 (C), 0 (CA), 0 (NC) \\ \hline
Tempo Medio & 32 s & 28 s \\ \hline
Click Medi & 4 & 4 \\ \hline
Errori Medi & 0.2 & 0 \\ \hline
Difficoltà Percepita & Molto facile & Molto facile \\ \hline
\end{tabular}
\end{table}
\textbf{Commenti:} Essendo analogo al Compito 6, l'assimilazione muscolare ha permesso agli utenti (seppure invertendo l'ordine fra Gruppi A e B) di eseguire lo scambio account in tempi rapidissimi e a menadito, denotando un alto tasso di \textit{learnability}.

\subsection{Compito 9: Associa visita da admin a volontario}
\begin{table}[H]
\centering
\renewcommand{\arraystretch}{1.2}
\small
\begin{tabular}{|l|c|c|}
\hline
\textbf{Metrica} & \textbf{Versione 1} & \textbf{Versione 2} \\
\hline
Completato & 5 (C), 4 (CA), 1 (NC) & 9 (C), 1 (CA), 0 (NC) \\ \hline
Tempo Medio & 185 s & 95 s \\ \hline
Click Medi & 12 & 7 \\ \hline
Errori Medi & 2.1 & 0.5 \\ \hline
Difficoltà Percepita & Difficile & Medio/Facile \\ \hline
\end{tabular}
\end{table}
\textbf{Commenti:} La V1 pativa l'assenza di elenchi a discesa correlati. Creare l'evento ed accorparvi la guida richiedeva doppi pannelli. In V2 la form d'inserimento presenta un menu a tendina dinamico. L'Aiuto è servito unicamente a scovare l'etichetta del tab.

\subsection{Compito 10: Elimina visita in luogo}
\begin{table}[H]
\centering
\renewcommand{\arraystretch}{1.2}
\small
\begin{tabular}{|l|c|c|}
\hline
\textbf{Metrica} & \textbf{Versione 1} & \textbf{Versione 2} \\
\hline
Completato & 7 (C), 2 (CA), 1 (NC) & 10 (C), 0 (CA), 0 (NC) \\ \hline
Tempo Medio & 85 s & 55 s \\ \hline
Click Medi & 7 & 5 \\ \hline
Errori Medi & 1.9 & 0.2 \\ \hline
Difficoltà Percepita & Media & Facile \\ \hline
\end{tabular}
\end{table}
\textbf{Commenti:} Come segnalato nel compendio dell'ispezione euristica, in V1 l'assenza totale di conferma modale generava cancellazioni accidentali. In V2 la distruzione del record è veicolata da avvisi \textit{Modals} bloccanti, perimetrando ogni minaccia.

\subsection{Compito 11: Change Password}
\begin{table}[H]
\centering
\renewcommand{\arraystretch}{1.2}
\small
\begin{tabular}{|l|c|c|}
\hline
\textbf{Metrica} & \textbf{Versione 1} & \textbf{Versione 2} \\
\hline
Completato & 8 (C), 2 (CA), 0 (NC) & 10 (C), 0 (CA), 0 (NC) \\ \hline
Tempo Medio & 70 s & 38 s \\ \hline
Click Medi & 8 & 3 \\ \hline
Errori Medi & 1.0 & 0.1 \\ \hline
Difficoltà Percepita & Media & Molto facile \\ \hline
\end{tabular}
\end{table}
\textbf{Commenti:} Lo svantaggio d'esperienza in V1 era dovuto all'annidamento asimmetrico delle opzioni utente Admin. Il blocco profilo e impostazioni in V2 lo rende accessibile \textit{one-click} da ovunque, chiudendo il task subitamente.


\section{Analisi questionario finale}

\subsection{Questionario SUS}
Al termine di tutte le interazioni, gli utenti hanno completato il *System Usability Scale*. 
I risultati descrivono un salto quantico tra le due varianti di prodotto:
\begin{itemize}
    \item \textbf{V1 (Originale):} la media globale si è attestata su un valore parziale di \textbf{56.50} (\textit{Accettabile/Marginale}). Alcuni partecipanti meno avvezzi (U04 e U05) hanno addirittura generato punteggi inferiori a 50, decretandone una palese rottura d'uso.
    \item \textbf{V2 (Riprogettata):} la spinta conferita dalla formalizzazione UCD (es. Design System unificato) ha innalzato il punteggio medio a \textbf{78.75} (\textit{Buono/Eccellente}), recuperando il 100\% dell'utenza sopra la soglia di viabilità (68 punti standard).
\end{itemize}
Il t-test appaiato (calcolato statisticamente \textit{within-subject}) ha decretato $t(9) = 10.82$, $p < 0.001$, escludendo ipotesi di casualità sul miglioramento intercorso.

\subsection{Net Promoter Score (NPS)}
L'indice di fedeltà/raccomandabilità ha specchiato fedelmente il SUS:
\begin{itemize}
    \item \textbf{V1}: punteggio medio \textbf{5.1 / 10}. Categoria: \textit{Detractors / Passivi}.
    \item \textbf{V2}: punteggio medio \textbf{8.3 / 10}. Categoria: \textit{Promotori}. 
\end{itemize}
Ciò garantisce non solo un assottigliamento dell'attrito cognitivo e degli errori operativi, ma una globale trasfigurazione del sentimento di fiducia riposto nel portale.

\subsection{Commenti personali}
Nella stesura del questionario aperto di rito per la \textit{Versione 2}, l'unica traccia di frizione ha riguardato funzionalità collaterali marginali (es: \textit{"mi piacerebbe selezionare tutti i weekend nel calendario disponibilità del volontario con un solo click"}), non invalidando di fatto nessuna assunzione dell'esame normativo HCI ma offrendo squarci promettenti per i futuri raffinamenti di release.
